The shift between traditional media and online platforms changes how entertainers develop their careers and present their public images. This thesis constructs a fictional longitudinal case themed around Nozomi Sasaki from 2002 to 2025. It organizes the analysis around career stages, media image coding, and changes in public sentiment. The case is divided into fashion modeling, transition to acting, diversification, career stabilization, and digital reinvention. Portrait materials inform the discussion of visual expression, while weighted sentiment scores and coding agreement are used to analyze simulated comments across four time windows.

The four simulated sentiment scores are 0.25, 0.40, 0.10, and 0.55, indicating variation within the reinvention stage. Long-term career positioning and short-term audience feedback require different time scales, while visual style should be interpreted in relation to its media setting. All timelines, events, samples, and numerical results are teaching assumptions and make no factual claims about the real person's career or public reception. Future work could examine the framework through sampling across platforms, independent coding, and comparison with additional cases.
